% Lösungen Einheit 1 von 4 – Länge, Masse, Geld (zusammenbau v0.1)
\einheitenkopf{Einheit 1 von 4 · Länge, Masse, Geld}

\erg{14}{a) $2$ m}
\erg{15}{a) Meter \quad b) $70$ dm \quad c) $3$ cm \quad d) $7{,}5$ m \quad e) $3\,000$ m \quad f) $5\,000$ kg \quad g) $345$ ct \quad h) $345$ cm \quad i) $235$ cm \quad j) $0{,}04$ m $= 4$ cm, also $0{,}04$ m $<$ $40$ cm \quad k) $55 \cdot 2{,}54$ cm $= 139{,}7$ cm $\approx 1{,}40$ m}
\erg{16}{a) Länge; unterstrichen: Eine Längeneinheit entspricht $30$ cm. \quad b) $6 \cdot 25 = 150$ cm \quad c) $7 - 2 = 5$ Längeneinheiten, $5 \cdot 4 = 20$ cm \quad d) $f(3) = -2{,}25 + 6 + 1{,}5 = 5{,}25$ m \quad e) Radius $f(2) = 2$, Durchmesser $4$ Längeneinheiten, $4 \cdot 3 = 12$ cm \quad f) Abstand $0{,}25$ Längeneinheiten, vorletzte Stange rechts bei $x = 0{,}5$: $f(0{,}5) = 1{,}6875$, Länge $1{,}6875 \cdot 1{,}6 = 2{,}7$ m; $2{,}4$ m Höhe sind $1{,}5$ Längeneinheiten: $0{,}6x^4 - 3x^2 + 0{,}9 = 0$, mit $z = x^2$: $z \approx 0{,}3206$ (die zweite Lösung liegt außerhalb), $x \approx \pm 0{,}5662$, Länge $2 \cdot 0{,}5662 \cdot 1{,}6 \approx 1{,}81$ m}
\erg{17}{a) mm; cm; dm; m; km}
\erg{18}{a) $30$ cm; $1{,}8$ m; $105$ m (Heft, Fahrrad, Fußballfeld)}
\erg{19}{a) $95$ cm; $1\,150$ mm; $1{,}2$ m}
\erg{20}{a) $32 \cdot 2{,}54 = 81{,}28 \approx 81{,}3$ cm}
\erg{21}{a) $5{,}25$ m}
\erg{22}{a) Tom hat malgenommen statt geteilt; von der kleinen in die große Einheit wird geteilt: $450 : 100 = 4{,}5$, also $450$ cm $= 4{,}5$ m \quad b) Jan hat die Längeneinheiten als Zentimeter übernommen und den Maßstab nicht angewendet; richtig: $7 \cdot 0{,}4 = 2{,}8$ cm \quad c) Die $9$ steht bei $0{,}09$ m an der Hundertstelstelle, das sind $9$ cm, nicht $90$ cm. \quad d) Eine Längeneinheit ist nur das Kästchenmaß; was sie in Wirklichkeit bedeutet, legt erst der Satz „eine Längeneinheit entspricht …“ im Text fest – bei jeder Aufgabe neu. \quad e) $128 + 2 \cdot 2 = 132$ cm; die Nische ist $135$ cm breit, Luft $135 - 132 = 3$ cm – das Regal passt \quad f) $2{,}06$ m; $2$ m $6$ cm}
